\documentclass{style}
\usepackage[utf8]{inputenc}
\usepackage[T1]{fontenc}
\usepackage{geometry}
\usepackage{graphicx}
\usepackage{float}
\usepackage{hyperref}
\usepackage[english,indonesian]{babel}

\title{Optimasi Pemrosesan Template LaTeX Menggunakan Pendekatan Berbasis AI}
\author{Academic Support Agent}
\date{\today}

\begin{document}
\maketitle

\begin{abstract}

\end{abstract}

\chapter{PENDAHULUAN}


\chapter{TINJAUAN PUSTAKA}


% Original content (commented out):
% \documentclass[12pt,a4paper]{style}
\input{preamble.tex} % all packages loaded and new definitions are in the preamble file

\title{\Large \bf
Interesting proposal title % proposal title goes here
}

\begin{document}

\section{\bf Goal of the experiment}
\label{sec:goal}
\textbf{The \code{lorem ipsum} text can be deleted by removing all \code{lipsum} commands.}
\lipsum[1-1]



\section{\bf Background}
\label{sec:background}

\begin{wrapfigure}{R}{0.55\textwidth}
\vspace{-5mm} % move picture up
 \includegraphics[width=0.55\textwidth]{./images/science.jpg}
 \vspace{-4mm} % spacing between figure and caption
\caption{The caption is in font size 10 to separate better from the text. The images can be loaded from the images folder. By choosing the \code{vspace} value (e.g.\qty{-2}{mm}), the picture can be moved up, the spacing between the image and the caption can be changed, and the space to the surrounding text below the caption can be adopted.
\label{fig:science}}
\vspace{0mm} % adjust space between caption and main text
\end{wrapfigure}

%Non-ablative ultrafast laser exposure can be used to manipulate a material's structure locally with incredible flexibility given that the laser parameters like focal spot size, pulse energy, polarization state and repetition rate can be tuned with ease using commercial lasers. 

\lipsum[1-2] And here is a reference to the figure with the glassware: c.f.~Fig.~\ref{fig:science}.


\section{\bf Experimental method; specific requirements}
\label{sec:method}
\lipsum[3]


\section{\bf Results expected and data management plan}
\label{sec:results}

\begin{figure}
 \begin{center}
 \includegraphics[width=0.8\textwidth]{./images/technology.jpg}
  \vspace{-2mm}
\caption{It is nicer to load pdf or other vector graphics, but jpeg format works as well. This is an example of an image that is not surrounded by text like the one in section \ref{sec:background}.
\label{fig:technology}}
 \end{center}
\vspace{-2mm}
\end{figure}

Here is an example on citing a reference in the text: \cite{sakai2000}. It is referenced to the bib file and shown in the list of references at the end of the document. And another reference can be found here \cite{diels2006}. \lipsum[4] \\

Also, a table is shown, c.f.~Tab.~\ref{table}. \lipsum[5] \\

\begin{table} [htbp]
\centering
    {\rowcolors{3}{white}{LightColor}
    \begin{tabularx}{0.4\textwidth}{X c}
    % one X is needed to allow for automatic stretching of the columns according to the defined textwidth
        \toprule
        \textbf{facility} & \textbf{opening year} \\
        \midrule
        SINQ & 1996 \\
        SLS & 2001 \\
        UCN & 2011 \\
        SwissFEL & 2018 \\
        \bottomrule
    \end{tabularx}
    \caption{Example table containing a few numbers of PSI facilities. Alternating the background colors for each row facilitates a clearer interpretation \label{table}}
    }
\end{table}
\lipsum[8-9]

\section{\bf Estimate and justification of the beamtime}
\label{sec:estimation}
\lipsum[5-6]

\footnotesize % bibliography in smaller fontsize
\bibliographystyle{unsrt}
%\usepackage[style=ieee]{citation-style-language}
\bibliography{references.bib} 
\normalsize


\end{document}


\rowcolor{LightColor} 

\end{document}